% ECONOMICS_PROBLEM: {"id": "C72-52", "title": "Games with constrained Nash equilibria for every feasible outside option", "jel_code": "C72", "source_id": "OP-0224"}
% !TeX program = xelatex
\documentclass[11pt,a4paper]{article}
\usepackage[margin=23mm]{geometry}
\usepackage{fontspec}
\setmainfont{Latin Modern Roman}
\usepackage{amsmath,amssymb,mathtools}
\usepackage{xcolor,enumitem,booktabs,longtable,array,xurl,hyperref}
\definecolor{muted}{HTML}{53636C}
\hypersetup{colorlinks=true,urlcolor=blue,linkcolor=blue}
\setlength{\parindent}{0pt}
\setlength{\parskip}{5pt}
\newcommand{\R}{\mathbb R}
\newcommand{\E}{\mathbb E}
\newcommand{\N}{\mathbb N}
\newcommand{\ind}{\mathbf 1}
\newcommand{\Var}{\operatorname{Var}}
\newcommand{\Cov}{\operatorname{Cov}}
\newcommand{\supp}{\operatorname{supp}}
\DeclareMathOperator*{\argmin}{arg\,min}
\DeclareMathOperator*{\argmax}{arg\,max}
\newcommand{\entrymeta}[3]{{\sffamily\small\color{muted}\textbf{\texttt{#1}}\quad|\quad #2\par #3\par}}
\newcommand{\Problemref}[1]{\texttt{#1}}
\newcommand{\JELref}[2]{\texttt{#2} (JEL \texttt{#1})}
\begin{document}
\section*{Games with constrained Nash equilibria for every feasible outside option}
\textbf{Problem ID:} \texttt{C72-52}\quad \textbf{JEL:} \texttt{C72}\par
% BEGIN ECONOMICS STATEMENT
\label{problem:OP-0224}
% GAME_THEORY_PROBLEM: {"local_id": "GT-094", "original_number": 94, "kind": "P", "entry_kind": "statement_only", "published": false, "review_required": true, "category": "game_theory"}
\label{game:GT-094}
{\sffamily\small\color{muted}
\noindent\textbf{ID:} GT-094. \textbf{Category:} Strategic matching and equilibrium existence.
\textbf{Status:} P; structural characterization requested in the inspected source.
\par}
\subsection*{Definitions and mathematical statement}
Let $X,Y$ be nonempty compact metric strategy spaces and let
$f,g:X\times Y\to\mathbb R$ be continuous payoffs. For outside-option payoffs $(a,b)\in\mathbb R^2$, put
\[
 F(a,b)=\{(x,y):f(x,y)\geq a,\ g(x,y)\geq b\}.
\]
A \emph{constrained Nash equilibrium} at $(a,b)$ is a pair $(x,y)\in F(a,b)$ such that
\begin{align*}
 f(x,y)&=\max\{f(z,y):z\in X,\ g(z,y)\geq b\},\\
 g(x,y)&=\max\{g(x,w):w\in Y,\ f(x,w)\geq a\}.
\end{align*}
The maxima exist at a feasible pair by compactness and continuity. Write $\operatorname{CNE}_G(a,b)$ for this set. A game is \emph{feasible}, in the source's terminology, if
\[
 \forall(a,b)\in\mathbb R^2,
 \qquad F(a,b)\neq\varnothing\Longrightarrow
                 \operatorname{CNE}_G(a,b)\neq\varnothing.
\]
Characterize these games by structural properties of $(X,Y,f,g)$, rather than by a restatement of the quantified equilibrium condition. A finite-dimensional specialization takes $X,Y$ to mixed-strategy simplexes of finite action sets and $f,g$ to bilinear expected payoffs; a characterization or recognition procedure for that subclass is a distinct, precise subproblem.

\subsection*{Short English statement}
Which bilateral games admit an equilibrium for every attainable pair of participation requirements?

\subsection*{Variants and scope boundaries}
Each unilateral deviation must preserve the \emph{other} player's outside-option requirement. This differs from an ordinary Nash equilibrium and from imposing only a player's own participation constraint. A feasible outside-option pair means joint attainability, not separate attainability in different profiles. Pure finite strategy spaces and mixed simplexes define different game classes. Characterization of a bilateral game is distinct from existence of a stable population matching built from many such games.

\subsection*{Source relationship and status}
Definitions 11--13 introduce constrained equilibria and feasible games, and the concluding discussion requests a characterization of feasible strategic games. The compact continuous formulation makes the maximization operations well-defined. The source's positive results for particular payoff classes are not a characterization of all games.

\subsection*{References}
\begingroup\footnotesize\raggedright
\noindent [S1] Felipe Garrido-Lucero and Rida Laraki, \emph{Stable Matching Games} (2020; inspected manuscript revision v4, 2025), arXiv:2008.01680v4. Inspected locators: Section 5.2, Definitions 11--12; Section 5.3, Definition 13; concluding discussion in Section 6.
\url{https://arxiv.org/pdf/2008.01680v4}
\par\endgroup

\subsection*{Common conventions: Source mathematical conventions}

\textbf{Finite games.} For a finite set $A$, $\Delta(A)$ is its probability
simplex. Players choose independent mixed strategies in Nash equilibrium;
correlation is allowed only when explicitly part of the solution concept.
In a game with payoff functions $u_i$ and strategy profile $x$, an additive
$\varepsilon$ Nash equilibrium satisfies
\[
 u_i(x)\geq u_i(x_i',x_{-i})-\varepsilon
 \quad\text{for every player $i$ and every admissible deviation $x_i'$}.
\]
For numerical approximation, payoffs are normalized as locally specified.
An $\varepsilon$ payoff residual is distinct from distance at most
$\varepsilon$ to the set of exact equilibria. Point convergence, convergence
of empirical play, and convergence in distance to an equilibrium set are
also distinct.

\textbf{Computational representations.} Unless stated otherwise, a complexity
question takes rational data with binary encoding; polynomial time is measured
in total bit length. A fixed-accuracy polynomial algorithm is not necessarily
polynomial in $1/\varepsilon$, and neither is necessarily polynomial in
$\log(1/\varepsilon)$. Succinct games and value, demand, or sampling oracles
must declare their interfaces. Exact real solutions require an explicit
representation; an unspecified list of arbitrary real numbers is not a
finite computational output. A claimed algorithmic barrier is meaningful
only for its specified model and reductions.

\textbf{Dynamic games.} Histories, information, observation, and allowable
randomization are part of the model. A stationary strategy depends only on
the current state; a positional strategy in a deterministic graph game
chooses one action at each controlled vertex. Nash equilibrium at an initial
state and equilibrium after every history are different requirements.
An expectation of a pathwise $\liminf$ payoff, a $\liminf$ of expected averages,
a limiting expected average, and a uniform finite-horizon equilibrium are
different criteria. None is silently substituted for another.

\textbf{Allocation and mechanisms.} A complete allocation assigns every item
exactly once unless otherwise stated. Zero-valued goods, negative values,
capacity constraints, feasibility, lotteries, and copies of goods affect
fairness definitions and are specified locally. Dominant-strategy incentive
compatibility, Bayesian incentive compatibility, universal truthfulness,
and truthfulness in expectation impose different inequalities. Whenever
an objective ratio has zero denominator, its convention is stated or those
instances are excluded.

\textbf{Infinite-dimensional models.} Expectations, integrals, stochastic
equations, and optimization objectives require their local regularity,
integrability, filtrations, and admissible domains. A backward--forward
mean-field system is a boundary-value problem; it is not an initial-value
semiflow without an additional construction. Long-time, large-population,
and vanishing-noise limits require an order of limits and a convergence
topology. If a minimum need not be attained, the objective is its infimum
or supremum and the approximation requirement is stated separately.
% END ECONOMICS STATEMENT
\end{document}
